\documentclass[12pt,a4paper,fontset=none]{ctexart}

\usepackage[a4paper,margin=2.35cm]{geometry}
\usepackage{amsmath}
\usepackage{amssymb}
\usepackage{booktabs}
\usepackage{longtable}
\usepackage{array}
\usepackage{enumitem}
\usepackage{hyperref}
\usepackage{xcolor}
\usepackage{listings}

\setCJKmainfont{Microsoft JhengHei}
\setCJKsansfont{Microsoft JhengHei}
\setCJKmonofont{Microsoft JhengHei}

\hypersetup{
  colorlinks=true,
  linkcolor=blue!55!black,
  urlcolor=blue!55!black
}

\lstdefinestyle{cmd}{
  basicstyle=\ttfamily\small,
  breaklines=true,
  frame=single,
  columns=fullflexible
}

\setlist[itemize]{leftmargin=2em}
\setlist[enumerate]{leftmargin=2em}

\title{\textbf{shogiAI 數據科學理論與應用報告}\\
\large Policy-only 將棋模型、資料前處理與 Alpha-beta 搜尋整合}
\author{shogiAI 專案}
\date{2026 年 6 月}

\begin{document}
\maketitle

\begin{abstract}
本報告說明 shogiAI 專案的數據科學流程與 AI 應用。專案以大量 CSA 將棋棋譜作為資料來源，從 MySQL 匯出「局面、下一手」訓練樣本，將局面編碼成 $43 \times 9 \times 9$ 張量，將走法編碼成 13,689 類 move label，訓練 policy-only 卷積神經網路。新版模型移除 value head，專注於預測實戰下一手與排序合法候選手；局面計算則交由 alpha-beta 搜尋、quiescence search、transposition table、killer move、history heuristic 與手工評估函數完成。最新遠端訓練使用 4,855,173 筆樣本，測試 Top-1 accuracy 為 41.10\%，Top-5 accuracy 為 74.49\%。
\end{abstract}

\tableofcontents
\newpage

\section{數據科學目標}

shogiAI 的數據科學目標是讓模型從人類實戰棋譜中學會「某個局面下，下一手最可能怎麼走」。新版模型採用 policy-only 設計，也就是模型只輸出走法機率，不再輸出局面勝率或 value。

模型可表示為：

\[
f_\theta(s) = p
\]

其中 $s$ 是目前局面，$p$ 是長度 13,689 的 policy logits。每一個 logit 對應一種可能走法。推論時，系統只取目前局面的合法手，讀取這些合法手對應的 policy 分數，再排序給 alpha-beta 搜尋使用。

\section{問題定義}

本專案的機器學習任務是多類別分類。

\begin{longtable}{p{0.28\textwidth}p{0.62\textwidth}}
\toprule
\textbf{項目} & \textbf{內容}\\
\midrule
\endhead
輸入 & 將棋局面，編碼為 $43 \times 9 \times 9$ 張量。\\
輸出 & 下一手 move label，總類別數為 13,689。\\
標籤來源 & 人類實戰棋譜中的下一手。\\
學習方法 & Supervised learning / imitation learning。\\
Loss & CrossEntropyLoss。\\
應用方式 & 排序合法手，輔助 alpha-beta 搜尋。\\
\bottomrule
\end{longtable}

這個任務不等於直接訓練一個能獨立下棋的模型。將棋局面往往有許多合理手，實戰下一手也不一定是唯一最佳手。因此模型的主要價值是把合理候選手排到前面，讓搜尋更有效率。

\section{資料來源與最新規模}

資料來源包含本地 CSA 檔案、遠端爬蟲蒐集的 CSA 檔案，以及 MySQL 中已匯入的棋譜資料。最新 policy-only 訓練是在遠端大數據主機上，從 MySQL 重新匯出資料集後執行。

\begin{longtable}{p{0.45\textwidth}p{0.40\textwidth}}
\toprule
\textbf{項目} & \textbf{最新數值}\\
\midrule
\endhead
匯出資料集 & \texttt{out/policy\_dataset.policy\_only\_latest.npz}\\
對局數 & 42,534 局\\
訓練樣本總數 & 4,855,173 筆\\
非法資料列 & 0 筆\\
\texttt{TORYO} 樣本 & 4,829,692 筆\\
\texttt{DRAW} 樣本 & 104 筆\\
\texttt{SENNICHITE} 樣本 & 25,377 筆\\
資料集大小 & 約 332 MB\\
\bottomrule
\end{longtable}

這些樣本是從資料庫的 \texttt{positions} 與 \texttt{moves} join 得到。每筆樣本代表「某一手前的局面」與「該局面下實戰走出的下一手」。

\section{資料前處理流程}

資料前處理主要由 \texttt{src/export\_policy\_dataset.py} 與 \texttt{src/csa\_preprocess.py} 負責。整體流程如下：

\begin{center}
\begin{tabular}{c}
MySQL \texttt{positions} + \texttt{moves} \\
$\downarrow$ \\
SFEN 重建 \texttt{cshogi.Board} \\
$\downarrow$ \\
合法手檢查 \\
$\downarrow$ \\
局面張量 \texttt{states} \\
$\downarrow$ \\
下一手 label \texttt{moves} \\
$\downarrow$ \\
壓縮保存為 \texttt{.npz}
\end{tabular}
\end{center}

核心 SQL 邏輯是將局面與下一手對齊：

\begin{lstlisting}[style=cmd]
SELECT p.game_id, p.move_number, p.sfen, m.usi_move, g.result
FROM positions p
JOIN moves m
  ON m.game_id = p.game_id
 AND m.move_number = p.move_number + 1
JOIN game_records g
  ON g.game_id = p.game_id
ORDER BY p.game_id, p.move_number;
\end{lstlisting}

在 Python 中，系統會用 \texttt{cshogi.Board(sfen)} 還原棋盤，用 \texttt{board.move\_from\_usi(usi\_move)} 取得 move，再以 \texttt{board.is\_legal(move)} 檢查合法性。非法資料不會進入訓練集。

\section{局面特徵編碼}

神經網路不能直接讀 SFEN 字串，因此系統將局面轉成固定大小張量：

\[
state \in \mathbb{R}^{43 \times 9 \times 9}
\]

43 個 planes 的意義如下：

\begin{longtable}{p{0.25\textwidth}p{0.60\textwidth}}
\toprule
\textbf{Plane 範圍} & \textbf{意義}\\
\midrule
\endhead
0--13 & 自方 14 種棋子在棋盤上的位置。\\
14--27 & 對方 14 種棋子在棋盤上的位置。\\
28--34 & 自方 7 種持駒數量。\\
35--41 & 對方 7 種持駒數量。\\
42 & 行棋方 plane。\\
\bottomrule
\end{longtable}

這種表示法類似影像中的多通道資料。棋盤是 $9 \times 9$ 空間，棋種、陣營、持駒與行棋方則作為 channel。CNN 可以在這種表示中學習局部形狀，例如攻防、升變區、玉周圍防守與棋子協同。

\section{行棋方視角旋轉}

專案預設 \texttt{orient\_to\_turn=True}。若輪到後手行棋，系統會把棋盤旋轉 180 度，讓模型永遠從「目前要下棋的一方」的視角看盤。

這個設計有三個好處：

\begin{enumerate}
  \item 減少先手與後手對稱造成的學習負擔。
  \item 讓同一戰型在不同顏色下能共享特徵。
  \item 讓模型輸出永遠表示「自己下一手」而不是固定先手或後手。
\end{enumerate}

例如後手局面會被轉成自方在下方的視角，後手的棋子會進入「自方棋子」planes。這使 policy network 不必另外學一套後手座標系。

\section{走法標籤設計}

Policy head 需要固定輸出維度，但每個局面的合法手數量不同。因此系統將所有可能走法編成固定 label，再在推論時只篩選合法手。

\subsection{一般移動}

一般移動由來源格、目標格與是否升變組成：

\[
81 \times 81 \times 2 = 13,122
\]

其中 81 是棋盤格數，乘以 2 是因為同一個來源與目標可能有升變與不升變兩種。

\subsection{打入持駒}

將棋有持駒打入規則。可打入棋子共有 7 種，目標格有 81 格：

\[
7 \times 81 = 567
\]

\subsection{總類別數}

\[
13,122 + 567 = 13,689
\]

因此模型最後一層輸出 13,689 維 logits。訓練時使用實戰下一手的 label 作為 target。

\section{從 Policy/Value 改成 Policy-only}

專案早期曾設計 policy/value network，模型同時輸出下一手 policy 與局面 value。Value head 的概念是用最終勝負作為標籤，讓模型學習目前行棋方最後會贏、輸或和棋：

\[
z \in \{-1,0,1\}
\]

後來實作與測試後，專案改成 policy-only。主要原因如下。

\subsection{Value 標籤訊號太粗}

Value label 來自一盤棋最後結果，但中盤局面距離終局很遠。某一方最後獲勝，不代表他在中盤每一手都真的佔優。用最終勝負硬套到所有中間局面，會讓 value target 含有大量噪音。

\subsection{投降資料使局面價值不精準}

許多棋譜以 \texttt{TORYO} 結束。投降能指出勝負方，但無法提供每個中間局面的精確勝率。尤其在將棋中，局面可能因一手戰術而突然逆轉，用整局結果回填所有局面會讓 value head 學到不穩定訊號。

\subsection{小型 CNN 容量有限}

目前模型是輕量 CNN，適合快速訓練與排序候選手，但要準確評估將棋複雜局面需要更深模型、更多搜尋標籤或自我對弈。若同時要求模型學 policy 與 value，可能分散有限模型容量。

\subsection{實際應用以候選手排序為主}

本專案 AI 最終仍使用 alpha-beta 搜尋計算變化。神經網路最需要做的是把好手排到前面，提高剪枝效率。局面評估則由傳統評估函數、搜尋深度與 quiescence search 負責。因此移除 value head 後，系統更簡單，也更符合目前需求。

\section{Policy-only 模型架構}

模型定義於 \texttt{src/policy\_model.py}，名稱為 \texttt{SmallPolicyNet}。架構如下：

\begin{lstlisting}[style=cmd]
Conv2d(43, 64, kernel_size=3, padding=1)
ReLU
Conv2d(64, 64, kernel_size=3, padding=1)
ReLU
Conv2d(64, 32, kernel_size=3, padding=1)
ReLU
Flatten
Linear(32 * 9 * 9, 512)
ReLU
Linear(512, 13689)
\end{lstlisting}

前半段卷積層保留棋盤空間結構，後半段 fully connected layer 將特徵轉成所有 move labels 的 logits。新版沒有 value head，因此模型只輸出 policy。

\section{訓練方法}

訓練腳本為 \texttt{src/train\_policy.py}。新版訓練流程只包含 policy 分類。

\subsection{Loss Function}

使用 cross entropy：

\[
\mathcal{L}_{policy}
=
-\log
\frac{\exp(\hat{p}_{a})}
{\sum_j \exp(\hat{p}_j)}
\]

其中 $a$ 是實戰下一手的 move label，$\hat{p}$ 是模型輸出的 logits。

\subsection{Optimizer}

使用 AdamW。AdamW 將 adaptive gradient 與 weight decay 分離，通常比傳統 Adam 加 L2 regularization 更穩定。

\subsection{Learning Rate Scheduler}

使用 CosineAnnealingLR。Learning rate 從初始值逐漸降低到最小值，使前期快速學習，後期微調模型。

\subsection{Game-level Split}

資料切分不是隨機切每一手，而是依 \texttt{game\_id} 切分。原因是同一盤棋相鄰局面非常相似，如果同一盤棋的前半放進訓練集、後半放進測試集，測試結果會過度樂觀。Game-level split 可以降低資料洩漏。

\section{最新訓練結果}

最新 policy-only 模型在遠端大數據主機訓練，並已同步回本地 \texttt{models/policy\_model.pt}。模型 SHA-256 為：

\begin{lstlisting}[style=cmd]
79bee4f2299667004803ec2265ba67a06c3236d08735372a7594a711b34b0ab0
\end{lstlisting}

\begin{longtable}{p{0.45\textwidth}p{0.40\textwidth}}
\toprule
\textbf{項目} & \textbf{數值}\\
\midrule
\endhead
Train samples & 3,882,230\\
Validation samples & 487,292\\
Test samples & 485,651\\
Test policy loss & 2.2751\\
Top-1 accuracy & 41.10\%\\
Top-3 accuracy & 65.24\%\\
Top-5 accuracy & 74.49\%\\
\bottomrule
\end{longtable}

初始局面的前五個 policy 候選手如下：

\begin{longtable}{p{0.30\textwidth}p{0.30\textwidth}}
\toprule
\textbf{走法} & \textbf{機率}\\
\midrule
\endhead
\texttt{7g7f} & 61.10\%\\
\texttt{2g2f} & 33.92\%\\
\texttt{5g5f} & 3.34\%\\
\texttt{2h7h} & 1.15\%\\
\texttt{9g9f} & 0.15\%\\
\bottomrule
\end{longtable}

\section{結果解讀}

Top-1 accuracy 表示模型直接猜中實戰下一手的比例。最新模型 Top-1 為 41.10\%，代表在測試局面中，模型第一候選手有約四成與實戰相同。

Top-5 accuracy 更適合評估本專案用途。將棋每個局面常有多個合理手，而模型主要用於搜尋前的 move ordering。Top-5 為 74.49\%，表示實戰下一手大約四分之三會出現在模型前五候選手中，這對 alpha-beta 剪枝很有幫助。

\section{模型推論與快取}

推論封裝為 \texttt{PolicyPredictor}。它負責：

\begin{itemize}
  \item 載入 \texttt{.pt} checkpoint。
  \item 將 \texttt{cshogi.Board} 轉成 $43 \times 9 \times 9$ state。
  \item 執行模型取得 logits。
  \item 只取合法手對應的 label 分數。
  \item 對合法手分數做 softmax，得到候選手機率。
  \item 依 score 排序回傳 \texttt{PolicyCandidate}。
  \item 用 \texttt{board.zobrist\_hash()} 做 prediction cache，減少重複局面推論。
\end{itemize}

這裡的 softmax 只在合法手集合內計算，因此機率表示「在目前合法手中，模型偏好此手的比例」。

\section{Alpha-beta 搜尋概論}

只靠 policy network 直接選最高機率手，容易漏看短期戰術，例如吃子交換、王手、詰將棋或防守手。因此 shogiAI 將神經網路與傳統搜尋結合。搜尋核心位於 \texttt{src/ai\_search.py}。

\subsection{Minimax 與 Negamax}

棋類遊戲是雙人零和遊戲。若一方得分越高，另一方就越低。Minimax 的想法是：我方選擇能讓自己最好的走法，對方也會選擇讓我最差的應手。

Negamax 是 minimax 的簡化形式，利用零和性質：

\[
score(s) = \max_a -score(T(s,a))
\]

其中 $T(s,a)$ 表示在局面 $s$ 下走 $a$ 後的新局面。每往下一層，視角會切換到對手，所以分數取負。

\subsection{Alpha-beta Pruning}

Alpha-beta 是 minimax 的剪枝最佳化。它維護兩個界線：

\begin{itemize}
  \item $\alpha$：目前已知我方至少能拿到的分數。
  \item $\beta$：目前對方允許我方拿到的上限。
\end{itemize}

若搜尋中發現某分支已不可能影響最終選擇，就停止展開該分支。剪枝不改變理論結果，但可大幅減少節點數。

\section{shogiAI 搜尋技術}

\subsection{Quiescence Search}

普通 alpha-beta 若剛好在戰術交換中停止，可能出現 horizon effect。例如葉節點看起來我方多一子，但下一手對方能吃回。Quiescence search 會在葉節點繼續延伸戰術手，例如吃子、升變或王手，直到局面較安靜再評估。

\subsection{Transposition Table}

不同走法順序可能到達同一局面。Transposition table 用 Zobrist hash 記錄已搜尋局面的分數、深度與最佳手。若搜尋再次遇到相同局面，可以直接使用快取結果或縮小 alpha-beta 範圍。

\subsection{Killer Move}

Killer move heuristic 記錄在同一搜尋深度曾造成 beta cutoff 的安靜手。若某手曾讓對手無法接受，之後同層搜尋可優先嘗試它，提高剪枝機會。

\subsection{History Heuristic}

History heuristic 統計某個走法在過去搜尋中造成 cutoff 的頻率與深度。越常在深層造成 cutoff 的走法，排序分數越高。

\subsection{Late Move Reduction}

若走法排序已經把較可能的好手放前面，後面排序很低且不是王手、不是戰術手的走法，可以先用較低深度搜尋。若低深度結果意外很好，再回頭完整搜尋。這可以節省大量時間。

\section{手工評估函數}

搜尋到葉節點時，需要評估局面分數。 \texttt{evaluate\_position} 不是只算子力，還加入多種將棋特徵：

\begin{itemize}
  \item \textbf{子力價值：} 飛、角、金、銀、桂、香、步與升變子的基礎價值。
  \item \textbf{持駒價值：} 將棋持駒可打入，因此持駒也有分數。
  \item \textbf{位置分數：} 棋子靠近中央、進入升變區或位置活躍會加分。
  \item \textbf{機動力：} 可行走法越多，局面越靈活。
  \item \textbf{玉安全：} 玉周圍保護、敵方壓力與直線攻擊。
  \item \textbf{懸空棋子：} 被攻擊但保護不足的棋子會扣分。
  \item \textbf{持駒壓力：} 持駒對敵玉周圍可打入威脅會加分。
  \item \textbf{節奏：} 輪到行棋方可能有主動權 bonus。
\end{itemize}

這些手工特徵提供基本棋力，使 AI 即使沒有神經網路，也能以 alpha-beta 下出合理走法。

\section{Policy Network 如何結合 Alpha-beta}

Policy network 不直接取代 alpha-beta，而是作為 move ordering。

\begin{center}
\begin{tabular}{c}
目前局面 \\
$\downarrow$ \\
列出合法手 \\
$\downarrow$ \\
Policy network 對合法手評分 \\
$\downarrow$ \\
合法手排序 \\
$\downarrow$ \\
Alpha-beta 依排序搜尋 \\
$\downarrow$ \\
選出實際走法
\end{tabular}
\end{center}

Alpha-beta 的效率非常依賴走法排序。如果好手越早被搜尋，$\alpha$ 或 $\beta$ 越快更新，就能越早剪掉其他分支。新版模型 Top-5 accuracy 達 74.49\%，表示它能有效把實戰合理手排在前面。

\section{Opening Book}

除了 policy network，AI 開局也可使用 opening book。Opening book 由 MySQL 統計前若干手常見局面與下一手頻率。若目前 SFEN 在 book 中，且某走法出現次數達到門檻，AI 可以直接採用該走法。

Opening book 的優點是：

\begin{itemize}
  \item 開局階段快速，不必消耗搜尋時間。
  \item 走法來自真實棋譜，較符合人類常見戰型。
  \item 可避免早期搜尋深度不足導致怪異開局。
\end{itemize}

\section{模型互打與評估}

\texttt{src/compare\_ai\_models.py} 可讓兩個模型互相對弈。它會交替先後手，避免先手優勢影響判斷。輸出包含：

\begin{itemize}
  \item 新模型勝場。
  \item 舊模型勝場。
  \item 和棋數。
  \item 新模型得分率。
  \item 平均手數。
  \item 每盤棋 USI 手順。
\end{itemize}

需要注意的是，Top-k accuracy 衡量的是「模仿實戰下一手」，模型互打才更接近棋力評估。但模型互打需要較多局數、固定深度與固定時間限制，結果才有統計意義。

\section{前端展示與應用}

前端 AI 對弈頁面會顯示：

\begin{itemize}
  \item 搜尋評分。
  \item 搜尋深度。
  \item 搜尋節點數。
  \item PV，也就是主要變化。
  \item Policy 候選手與機率。
\end{itemize}

新版前端已移除「局面值」欄位，因為模型不再輸出 value。這能避免使用者誤以為神經網路仍在估計勝率。

\section{限制與改進方向}

\subsection{模型容量限制}

目前模型是小型 CNN，訓練與推論速度快，但容量有限。若要提升 policy 準確率，可改用 residual CNN 或 attention-based 架構。

\subsection{資料標籤限制}

模型學的是人類實戰下一手，而不是數學上的最佳手。同一局面可能有多個好手，但資料只標一手。因此 Top-1 不應被解讀為唯一棋力指標。

\subsection{Legal Move Mask}

目前模型輸出所有 13,689 類，推論時再挑合法手。未來可在 loss 中加入合法手 mask 或改良 label 設計，讓模型更專注於實際可行走法。

\subsection{自我對弈}

若要讓模型不只模仿人類，可加入 self-play。透過目前 AI 搜尋產生新棋譜，再用結果訓練模型，逐步提升棋力。

\subsection{Value 的未來可能}

Value head 目前被移除，不代表永遠不能使用。若未來要重新嘗試 value，建議不要直接用最終勝負套到所有中間局面，而是使用搜尋評分、自我對弈結果、或更細緻的 bootstrap target。這樣 value 訊號會比單純終局勝負更穩定。

\section{結論}

shogiAI 的數據科學流程已從早期 policy/value 模型調整為 policy-only 模型。新版設計讓神經網路專注於從大量實戰棋譜中學習下一手分布，並把學到的知識用於 alpha-beta 搜尋的走法排序。最新資料集包含 4,855,173 筆樣本，測試 Top-1 accuracy 達 41.10\%，Top-5 accuracy 達 74.49\%。整體系統結合了資料庫、特徵工程、CNN 分類模型、opening book 與傳統棋類搜尋，使專案具備從資料蒐集到實際 AI 對弈的完整流程。

\appendix
\section{常用指令}

\subsection{匯出資料集}

\begin{lstlisting}[style=cmd]
python src\export_policy_dataset.py ^
  --output out\policy_dataset.policy_only_latest.npz ^
  --host 140.135.65.53 ^
  --port 3306 ^
  --user 11211213 ^
  --database DB11211213
\end{lstlisting}

\subsection{訓練 policy-only 模型}

\begin{lstlisting}[style=cmd]
python src\train_policy.py ^
  --input out\policy_dataset.policy_only_latest.npz ^
  --output out\policy_model.policy_only.pt ^
  --batch-size 256
\end{lstlisting}

\subsection{啟動 AI 瀏覽器}

\begin{lstlisting}[style=cmd]
python src\csa_browser_api.py ^
  --source mysql ^
  --db-host 140.135.65.53 ^
  --db-port 3306 ^
  --db-user 11211213 ^
  --db-name DB11211213 ^
  --policy-model models\policy_model.pt ^
  --policy-order-ply 2
\end{lstlisting}

\subsection{模型互打}

\begin{lstlisting}[style=cmd]
python src\compare_ai_models.py ^
  --old-model out\policy_model.prev.pt ^
  --new-model models\policy_model.pt ^
  --games 20 ^
  --depth 3 ^
  --time-limit-ms 1000
\end{lstlisting}

\end{document}
